\input{../tex-inputs/latex-leseansicht-vorspann}

\section[Arthur und Olga Schnitzler und Paul Marx an Gustav Schwarzkopf, 7. 9. 1913]{L04509 Arthur und Olga Schnitzler und Paul Marx an Gustav Schwarzkopf, 7. 9. 1913}
\nopagebreak\mylabel{L04509v}
\rehead{ }\normalsize\beginnumbering\briefempfaengerindex{Schwarzkopf, Gustav@\textsc{Schwarzkopf, Gustav}!zzzMarx, Paul@\emph{von Paul Marx}!1913-09-071@{7. 9. 1913}|(be}\briefempfaengerindex{Schwarzkopf, Gustav@\textsc{Schwarzkopf, Gustav}!zzzSchnitzler, Olga@\emph{von Olga Schnitzler}!1913-09-071@{7. 9. 1913}|(be}\briefempfaengerindex{Schwarzkopf, Gustav@\textsc{Schwarzkopf, Gustav}!zzzSchnitzler, Arthur@\emph{von Arthur Schnitzler}!1913-09-071@{7. 9. 1913}|(be}
\toendnotes[C]{\smallbreak\pagebreak[2]}
\correspDesc{Versand  durch Arthur Schnitzler, Olga Schnitzler, Paul Marx am 7. 9. 1913 in München
\newline{}Übermittlung  am 7. 9. 1913 in München
\newline{}Erhalt  durch Gustav Schwarzkopf im Zeitraum [8. 9. 1913 – 12. 9. 1913?] in Wien}\toendnotes[C]{\smallbreak}
\Standort{DLA, A:Schnitzler, HS.1985.1.1897.}
\physDesc{Bildpostkarte, 197 Zeichen
\newline{}Handschrift Arthur Schnitzler: Bleistift, deutsche Kurrent
\newline{}Handschrift Olga Schnitzler: Bleistift, lateinische Kurrent
\newline{}Handschrift Paul Marx: Bleistift, deutsche Kurrent
\newline{}Versand: Stempel: »\nobreak{}\oindex{München@\textbf{München}|pwk}München, 7. 9. 13., 5–6 N\nobreak{}«.  }
\pstart
           \noindent{}{\pb}\textcolor{gray}{\textbf{München\oindex{München@\textbf{München}|pw}\hspace*{2.5em}Reichenbachbrücke\oindex{Reichenbachbrücke@\textbf{Reichenbachbrücke}|pw}, Blick gegen die Au\oindex{Au@\textbf{Au}|pw}}}\pend
           \pstart{}{\pb}Herrn \textsc{Gustav Schwarzkopf}\pend{}\pstart{}Wien I\oindex{I., Innere Stadt@\textbf{I., Innere Stadt}|pw}\pend{}\pstart{}\textsc{Tiefer Graben 17}\oindex{Wien@\textbf{Wien}!I., Innere Stadt@\textbf{I., Innere Stadt}!Tiefer Graben 17@\textbf{Tiefer Graben 17}|pw}\pend{}{\bigskip}\vspace{1em}
\pstart
           \centering{}\textsc{München}\oindex{München@\textbf{München}|pw}{ }7/9 1913\pend
           \vspace{0.5em}
\pstart
           Herzliche Grüße u auf baldgs
      Wiederſehn!\pend
           \pstart Ihr \spacefill\mbox{Arthur}\pend{}\selectlanguage{ngerman}\vspace{1em}
\pstart
           \noindent{}{[}hs. O. Schnitzler:{]} Herzlichst Ihre \spacefill\mbox{Olga}.\pend
           \selectlanguage{ngerman}\vspace{1em}
\pstart
           \noindent{}{[}hs. Marx:{]} Herzlichſte Empfehlungen. Wie geht es Ihnen geſundheitlich? Ihr{\\}\spacefill\mbox{Marx}\pend
           \selectlanguage{ngerman}\endnumbering\briefempfaengerindex{Schwarzkopf, Gustav@\textsc{Schwarzkopf, Gustav}!zzzMarx, Paul@\emph{von Paul Marx}!1913-09-071@{7. 9. 1913}|)be}\briefempfaengerindex{Schwarzkopf, Gustav@\textsc{Schwarzkopf, Gustav}!zzzSchnitzler, Olga@\emph{von Olga Schnitzler}!1913-09-071@{7. 9. 1913}|)be}\briefempfaengerindex{Schwarzkopf, Gustav@\textsc{Schwarzkopf, Gustav}!zzzSchnitzler, Arthur@\emph{von Arthur Schnitzler}!1913-09-071@{7. 9. 1913}|)be}\mylabel{L04509h}
\begin{anhang}
\end{anhang}\newcommand{\dateiname}{L04509}\newcommand{\titel}{Arthur und Olga Schnitzler und Paul Marx an Gustav Schwarzkopf, 7. 9. 1913}\newcommand{\editorInnen}{Herausgegeben von Jahnke, SelmaMüller, Martin Anton}\input{../tex-inputs/latex-leseansicht-abspann}
